% ============================================================================
% 7. CONCLUSION — Journal version (expanded with future work)
% ============================================================================
\section{Conclusion and Future Work}
\label{sec:conclusion}

We presented DW-Bench, a two-tier benchmark for evaluating LLMs on graph
topology reasoning over data warehouse schemas.  Across six baselines and
three frontier models, four findings consistently emerge:

\begin{enumerate}
  \item \textbf{Tools dominate static contexts} for schema-level
    reasoning (89--90\% vs.\ 69--77\% micro-EM), confirming that
    interactive graph algorithm access is essential.  Static context
    injection cannot substitute for on-demand traversal, regardless
    of context length.

  \item \textbf{Static baselines exhibit a Triviality Illusion}: easy
    subtypes like \texttt{join\_path} inflate micro-EM scores, masking
    failures on compositional reasoning.  Practitioner reporting should
    always disaggregate by subtype difficulty.

  \item \textbf{A hard-question ceiling persists} at $\sim$60\% EM even
    with tools, compared to $>$99\% Oracle, identifying multi-hop
    compositional reasoning as the core bottleneck.  Budget ablation
    (ReAct-Code at 5 rounds) confirms the bottleneck is reasoning
    capacity, not interaction budget.

  \item \textbf{Tools are necessary but insufficient} for value-level
    questions: essential for data-access subtypes (100\% vs.\ 0\%) but
    provide no advantage on multi-hop structural reasoning (27\%).
    This points toward learned structural representations: GNN
    encoders or graph transformers: as the missing inductive bias.
\end{enumerate}

\subsection{Future Directions}

\textbf{GNN-Augmented Baselines.}
The reasoning-limited subtypes suggest that learned graph representations
(GNN encoders, graph transformers) could provide the structural inductive
bias that prompt-based approaches lack.  Training a GNN on schema graph
features and injecting learned embeddings as context is a natural
extension that DW-Bench provides the infrastructure to evaluate.

\textbf{Agentic Methods.}
Advanced agentic frameworks (ToG~\cite{sun2024tog},
GraphCoT~\cite{jin2024graphcot}) that decompose complex queries into
multi-step reasoning chains may close the gap on compositional subtypes
like \texttt{combined\_impact}; DW-Bench's per-subtype diagnostics
provide the necessary granularity to measure progress.

\textbf{Broader Model and Dataset Coverage.}
Extending evaluation to GPT-4 class models, open-weight LLMs (Llama,
Mistral), and additional real-world schemas with row-level lineage would
strengthen generality claims and enable Tier~2 evaluation beyond
synthetic data.

\textbf{Schema Evolution.}
Real data warehouses change continuously.  Evaluating whether LLMs
can reason about schema topology \emph{under perturbation} (table
additions, FK modifications, lineage updates) would extend DW-Bench
toward production-relevant settings.

\subsection{Reproducibility}

All datasets, generation scripts, evaluation code, baseline
implementations, and per-question results are publicly available at
\url{https://github.com/AJamal27891/dw-bench}.  Benchmark data will
also be released on HuggingFace at
\url{https://huggingface.co/datasets/AJamal27891/dw-bench}.
Results are deterministically reproducible
using greedy decoding with the specified model versions.

\paragraph{Dataset licenses.}
TPC-DS and TPC-DI are available under the TPC EULA (free for research).
OMOP CDM uses the Apache~2.0 license.  AdventureWorks is released under
the Microsoft Public License.  Syn-Logistics is our creation, released
under MIT.
